\chapter{Generative Adversarial Networks}

\section{Adversarial Formulation}

Generative adversarial networks (GANs) provide an alternative approach to generative modeling based on a game between two models.

\medskip

\noindent
\textbf{Core idea.}  
Two networks are trained simultaneously:
\begin{itemize}
    \item \textbf{Generator} $G_\theta$: maps noise $z \sim p(z)$ to data space
    \item \textbf{Discriminator} $D_\phi$: distinguishes real data from generated data
\end{itemize}

\medskip

\noindent
\textbf{Minimax objective.}
\[
\min_\theta \max_\phi \;
\mathbb{E}_{x \sim p_{\text{data}}}[\log D_\phi(x)]
+
\mathbb{E}_{z \sim p(z)}[\log(1 - D_\phi(G_\theta(z)))].
\]

\medskip

\noindent
\textbf{Interpretation.}
\begin{itemize}
    \item The generator tries to fool the discriminator
    \item The discriminator tries to distinguish real and fake samples
\end{itemize}

\medskip

\noindent
\textbf{Equilibrium.}  
At optimality, the generator produces samples indistinguishable from real data.

\medskip

\noindent
\textbf{Implicit modeling.}  
GANs define a sampling procedure without explicitly modeling $p_\theta(x)$.

\medskip

\noindent
\textbf{Insight.}  
Learning is framed as a game between two competing objectives.

\section{Divergence Minimization}

GANs can be understood as minimizing a divergence between distributions.

\medskip

\noindent
\textbf{Optimal discriminator.}  
For a fixed generator, the optimal discriminator is:
\[
D^*(x) = \frac{p_{\text{data}}(x)}{p_{\text{data}}(x) + p_\theta(x)}.
\]

\medskip

\noindent
\textbf{Resulting objective.}  
Substituting $D^*$ yields:
\[
\min_\theta \; \mathrm{JS}(p_{\text{data}} \,\|\, p_\theta),
\]
where $\mathrm{JS}$ denotes the Jensen--Shannon divergence.

\medskip

\noindent
\textbf{Interpretation.}
\begin{itemize}
    \item GAN training implicitly minimizes a divergence between distributions
    \item The generator improves as distributions become closer
\end{itemize}

\medskip

\noindent
\textbf{Limitations.}
\begin{itemize}
    \item Divergences may behave poorly when supports do not overlap
\end{itemize}

\medskip

\noindent
\textbf{Insight.}  
GANs connect game-theoretic learning with statistical divergence minimization.

\section{Training Instability}

Training GANs is often challenging due to the adversarial nature of the objective.

\medskip

\noindent
\textbf{Non-convex--non-concave optimization.}  
The minimax problem does not generally admit stable convergence.

\medskip

\noindent
\textbf{Mode collapse.}
\begin{itemize}
    \item Generator produces limited variety of outputs
    \item Fails to capture full data distribution
\end{itemize}

\medskip

\noindent
\textbf{Oscillatory dynamics.}
\begin{itemize}
    \item Generator and discriminator may fail to converge
    \item Parameters can cycle or diverge
\end{itemize}

\medskip

\noindent
\textbf{Gradient issues.}
\begin{itemize}
    \item Vanishing gradients when discriminator is too strong
    \item Unstable updates due to adversarial feedback
\end{itemize}

\medskip

\noindent
\textbf{Practical techniques.}
\begin{itemize}
    \item Careful balancing of generator and discriminator updates
    \item Label smoothing
    \item Gradient penalties
\end{itemize}

\medskip

\noindent
\textbf{Insight.}  
Adversarial training introduces dynamics that are fundamentally different from standard optimization.

\section{Variants and Improvements}

Many variants of GANs have been proposed to address training challenges and improve performance.

\medskip

\noindent
\textbf{Wasserstein GAN (WGAN).}
\begin{itemize}
    \item Replaces JS divergence with Wasserstein distance
    \item Improves stability and gradient behavior
\end{itemize}

\medskip

\noindent
\textbf{WGAN with gradient penalty.}
\begin{itemize}
    \item Enforces Lipschitz constraint via gradient penalties
\end{itemize}

\medskip

\noindent
\textbf{Conditional GANs.}
\begin{itemize}
    \item Generate samples conditioned on labels or inputs
\end{itemize}

\medskip

\noindent
\textbf{CycleGAN and image translation.}
\begin{itemize}
    \item Learn mappings between domains without paired data
\end{itemize}

\medskip

\noindent
\textbf{Style-based GANs.}
\begin{itemize}
    \item Introduce structured latent spaces
    \item Enable fine control over generated images
\end{itemize}

\medskip

\noindent
\textbf{Insight.}  
Improving GANs often involves modifying the objective, architecture, or training procedure.

\section{A Unifying Perspective}

GANs provide a distinct approach to generative modeling:

\begin{itemize}
    \item \textbf{Implicit models:} define sampling without explicit density
    \item \textbf{Adversarial training:} two networks in competition
    \item \textbf{Divergence minimization:} alignment of distributions
\end{itemize}

\medskip

\noindent
\textbf{Core components.}
\begin{itemize}
    \item Generator $G_\theta$
    \item Discriminator $D_\phi$
    \item Minimax objective
\end{itemize}

\medskip

\noindent
\textbf{Key insight.}  
Adversarial learning provides a powerful but challenging framework for generative modeling.

\section{Outlook}

This chapter introduced generative adversarial networks:
\begin{itemize}
    \item adversarial formulation,
    \item divergence-based interpretation,
    \item training challenges,
    \item key variants and improvements.
\end{itemize}

\medskip

In the next chapter, we study diffusion models, which provide a fundamentally different approach to generative modeling based on stochastic processes.

\medskip

\noindent
\textbf{Guiding principle.}  
Different generative frameworks offer different tradeoffs between tractability, stability, and sample quality.